\makeatletter\def\input@path{{../template/vsp-beamer}{../template/minimus}}\makeatother
\documentclass[red,tutorial,sectionpages,shtu]{vsp-beamer}
\graphicspath{{../template/vsp-beamer/assets/}}
\usepackage{minimus-section}
\usepackage{minimus-text}
\usepackage{minimus-math}
\usepackage{minimus-float}
\usepackage{minimus-code}
\usepackage{minimus-colorbox}
\usepackage{minimus-reference}
\usepackage{tikz}
\usetikzlibrary{arrows.meta,positioning,calc}
\lstset{basicstyle=\ttfamily\footnotesize,numberstyle=\tiny\color{gray},breaklines=true,xleftmargin=8pt}
\title{CS100 Recitation 5}
\subtitle{C语言内存管理：生命周期、指针与动态分配}
\VSPsetspeaker{Chaofan Li}{CS100 Teaching Assistant}
\VSPsetupfootline{CS100 Recitation 5}{Fall 2026}{Chaofan Li}
\begin{document}
\VSPtitleframe
\begin{frame}{今天我们要回答什么？}
    \begin{enumerate}
        \item 我们为什么需要指针？
        \item 运行时才知道需要多少空间，怎么办？
        \item 这些对象通常放在哪些memory区域？
        \item 一块动态分配的内存，最后应该由谁释放？
    \end{enumerate}
    \medskip
    \alert{看到一个指针，先问：它指向谁？那个对象还活着吗？}
\end{frame}
\begin{frame}{目录}
    \tableofcontents[hideallsubsections]
\end{frame}
\section{先看生命周期，再看内存布局}

\begin{frame}{先猜输出：每次进入函数，变量都会重新开始吗？}
\begin{Code}*{c}\lstinputlisting{code/lifetime.c}\end{Code}
\note{先留一分钟，让同学写下两行输出；不要先把static等同于全局可见。}
\end{frame}

\begin{frame}{答案：作用域与生命周期是两件事}
\begin{itemize}
\item 第一次：\texttt{1 1 11}；第二次：\texttt{1 2 12}。
\item \texttt{local}：每次进入函数体，开始新的automatic对象生命周期。
\item \texttt{count}：具有static storage duration，整个程序执行期间都存在。
\item \texttt{global}：同样具有static storage duration，值不会在函数返回时重置。
\end{itemize}
\medskip\alert{count只能在函数体内用名字访问，但它不会随着函数返回而消失。}
\end{frame}

\begin{frame}{Scope、storage duration、lifetime}
\begin{itemize}
\item \textbf{Scope（作用域）}：源码中，哪些位置能用这个名字？
\item \textbf{Storage duration（存储期）}：对象的存储存在多长时间？
\item \textbf{Lifetime（生命周期）}：对象在执行中何时存在、何时可以访问？
\end{itemize}
\medskip
\begin{tabular}{lll}
\textbf{对象} & \textbf{典型存储期} & \textbf{何时结束}\\[4pt]
普通块内局部变量 & automatic & 离开对应的块\\
全局变量／\texttt{static}局部变量 & static & 程序结束\\
\texttt{malloc}分配的对象 & allocated & 被释放或被成功重新分配
\end{tabular}

\note{本节限定普通非VLA示例，不引入thread storage duration。指针也是对象，拥有自己的生命周期。}
\end{frame}

\begin{frame}{Memory layout：常见Linux进程的示意图}

\begin{columns}[c]
\begin{column}{0.43\textwidth}
\centering
\begin{tikzpicture}[every node/.style={draw=VSPPrimary,minimum width=4.7cm,minimum height=0.60cm,align=center,font=\small},node distance=0.06cm]
\node (stack) {Stack：调用栈、部分局部对象};
\node[below=of stack] (gap) {\(\downarrow\)\quad 空间／其他映射\quad\(\uparrow\)};
\node[below=of gap] (heap) {Heap：动态分配区域};
\node[below=of heap] (bss) {.bss：零初始化静态对象};
\node[below=of bss] (data) {.data：非零初始化静态对象};
\node[below=of data] (ro) {.rodata：字符串字面量等只读数据};
\node[below=of ro] (text) {.text：机器指令};
\end{tikzpicture}
\end{column}
\begin{column}{0.52\textwidth}
高地址在上，低地址在下。\medskip
\begin{itemize}
\item 这是理解常见实现的模型。
\item 实际进程还包括共享库、内存映射等区域。
\item 动态分配也可能通过独立映射实现。
\item \alert{C标准不保证这些分区、顺序或增长方向。}
\end{itemize}
\end{column}
\end{columns}

\end{frame}

\begin{frame}{对象放在哪里？先说规则，再说常见实现}
\begin{Code}*{c}\lstinputlisting{code/locations.inc}\end{Code}\begin{itemize}
\item \texttt{g}通常在.data；\texttt{z}、\texttt{count}通常在.bss。
\item \texttt{local}和指针变量\texttt{p}通常在stack，也可能被优化到寄存器。
\item \alert{p与*p不是同一个对象。}动态分配的空间通常来自heap。
\end{itemize}

\note{片段省略stdlib.h与失败处理，只用于标识位置；后续完整示例会检查NULL。}
\end{frame}

\begin{frame}{.bss与.data：为什么两个程序大小差这么多？}
\begin{columns}[T]\begin{column}{0.53\textwidth}\begin{Code}*{c}\lstinputlisting{code/bss.c}\end{Code}\begin{Code}*{c}\lstinputlisting{code/data.c}\end{Code}\end{column}\begin{column}{0.43\textwidth}\begin{itemize}
\item 两组数组在运行时都需要存储。
\item 第一组默认全为0；第二组前三项非零，其余仍为0。
\item .bss通常只记录需要多少零初始化空间。
\item .data需要在文件中保存初始数据。
\end{itemize}
\end{column}\end{columns}\medskip\alert{可执行文件小，不代表程序运行时不占内存。}
\end{frame}

\begin{frame}{现场看一眼：size与objdump}
在\texttt{r5/}目录执行：\begin{Code}*{bash}\lstinputlisting{code/layout.sh}\end{Code}\begin{itemize}
\item \texttt{size}：比较text、data、bss所占空间。
\item \texttt{objdump -t}：观察符号\texttt{array}属于哪个section。
\item \texttt{int a;}与\texttt{int b = 0;}通常都进入.bss。
\end{itemize}

\note{工具演示约3分钟，只观察section归属。零初始化不是未初始化；具体文件大小取决于平台与工具链。}
\end{frame}

\begin{frame}{Stack frame：函数返回后发生了什么？}

\begin{columns}[c]
\begin{column}{0.43\textwidth}
\begin{tikzpicture}[node distance=0.12cm,every node/.style={draw=VSPPrimary,minimum width=4.6cm,minimum height=0.85cm,align=center}]
\node (c) {\texttt{function\_c}的frame};
\node[below=of c] (b) {\texttt{function\_b}的frame};
\node[below=of b] {\texttt{main}的frame};
\end{tikzpicture}
\end{column}
\begin{column}{0.52\textwidth}
\begin{itemize}
\item 函数调用通常建立stack frame（栈帧）。
\item 栈帧可能保存局部变量、返回地址与部分参数。
\item 参数也可能通过寄存器传递。
\item 函数返回后，普通局部对象生命周期结束。
\end{itemize}
\end{column}
\end{columns}
\medskip\alert{内存里的旧字节暂时没被覆盖，也不代表原来的对象还活着。}

\end{frame}

\begin{frame}{这段代码为什么不能返回局部变量的地址？}
\begin{Code}*{c}\lstinputlisting{code/return-local.inc}\end{Code}\begin{itemize}
\item 函数返回后，\texttt{x}的生命周期已经结束。
\item 调用者不能通过返回值继续访问\texttt{x}。
\item 若只需要数值，直接返回\texttt{int}；若需要长期存在的对象，重新设计存储与所有权。
\end{itemize}
\medskip\alert{不能根据某一次运行“恰好输出42”判断代码正确。}
\end{frame}

\section{运行时需要空间：动态分配}

\begin{frame}{为什么需要Dynamic allocation？}
\begin{itemize}
\item 输入告诉你有多少个元素，编译时还不知道。
\item 对象需要跨越一次函数调用继续存在。
\item 数组需要扩容，或者需要逐个建立链表节点。
\end{itemize}
\medskip
\begin{tabular}{lll}
 & \textbf{automatic局部对象} & \textbf{动态分配对象}\\[4pt]
创建 & 进入对应块时自动管理 & 调用\texttt{malloc/calloc}等\\
结束 & 离开对应块 & 调用\texttt{free/realloc}等\\
责任 & 注意作用域与生命周期 & 额外设计释放责任
\end{tabular}
\medskip\par\alert{选择依据是大小和生命周期需求，不是“heap一定比stack好”。}

\end{frame}

\begin{frame}{四个接口，先记住各自做什么}
\begin{Code}*{c}\lstinputlisting{code/api.inc}\end{Code}\begin{itemize}
\item \texttt{size\_t}：用于表示大小的无符号整数类型。
\item \texttt{malloc}：申请指定字节数；\texttt{calloc}：申请数组并将所有位清零。
\item \texttt{realloc}：调整已有分配；\texttt{free}：结束该次分配。
\item 本节示例均使用\alert{正数大小}，避免零大小分配的特殊规则。
\end{itemize}

\end{frame}

\begin{frame}{malloc：空间申请到了，不代表已经有初值}
\begin{Code}*{c}\lstinputlisting{code/malloc.c}\end{Code}\alert{检查NULL之后再访问；读取之前先写入有效值。}
\end{frame}

\begin{frame}{sizeof *p：申请的是对象的大小}
\begin{Code}*{c}\lstinputlisting{code/sizes.inc}\end{Code}\begin{itemize}
\item \texttt{sizeof p}是指针本身的大小，不是一个\texttt{int}的大小。
\item \texttt{sizeof *p}取所指类型的大小，修改指针类型时不必同步修改类型名。
\item 对这里的非VLA类型，\texttt{sizeof *p}不求值\texttt{*p}，不会解引用未初始化的\texttt{p}。
\item C中\texttt{void *}可转换成对象指针，不需要给\texttt{malloc}返回值强制类型转换。
\end{itemize}

\end{frame}

\begin{frame}{输入决定数组大小：乘法之前也要检查}
\begin{Code}*{c}\lstinputlisting{code/array.c}\end{Code}
\note{课堂输入协议为可由size_t表示的非负整数。边界检查防止字节数乘法回绕，不是完整的任意文本数字解析器。}
\end{frame}

\begin{frame}{calloc：按元素数量申请，并把所有位清零}
\begin{Code}*{c}\lstinputlisting{code/calloc.c}\end{Code}\begin{itemize}
\item 这里的\texttt{int}数组初始值全是0，输出\texttt{0 0}。
\item 所有位为0，不可推广成所有类型的语义零值；例如不保证得到空指针表示。
\end{itemize}

\end{frame}

\begin{frame}{realloc：扩大数组，原有数据会怎样？}
\begin{itemize}
\item \texttt{realloc(p, new\_size)}调整由分配函数返回、尚未释放的块。
\item 成功：保留新旧大小较小范围内的数据；新增部分未初始化。
\item 成功后原分配结束，\alert{旧指针及指向旧块的别名都不能继续使用}。
\item 失败且请求大小非零：返回\texttt{NULL}，原分配仍然有效。
\end{itemize}
\medskip\alert{它可能原地扩容，也可能搬家。代码必须同时支持两种情况。}
\end{frame}

\begin{frame}{为什么不能直接p = realloc(p, ...)？}
\begin{Code}*{c}\lstinputlisting{code/realloc-wrong.inc}\end{Code}\begin{itemize}
\item 失败时返回\texttt{NULL}，赋值覆盖了原来的指针。
\item 原块没有被释放，却可能失去唯一可用的地址，造成memory leak。
\item 用临时指针接住返回值，成功后再更新\texttt{p}。
\end{itemize}

\end{frame}

\begin{frame}{安全扩容：先接住，再更新}
\begin{columns}[T]
\begin{column}{0.48\textwidth}
\begin{Code}*{c}\lstinputlisting[firstline=1,lastline=10]{code/realloc.c}\end{Code}
\end{column}
\begin{column}{0.48\textwidth}
\begin{Code}*{c}\lstinputlisting[firstline=11,lastline=20,firstnumber=11]{code/realloc.c}\end{Code}
\end{column}
\end{columns}
\medskip
\alert{失败：原块仍有效，清理后退出。成功：更新指针，初始化新增元素。}
\note{固定大小不会溢出。提醒同学：不要在成功后比较旧p或沿用旧块内部指针；这版例子没有建立别名。}
\end{frame}

\begin{frame}{选读：分配器大概在做什么？}
\begin{itemize}
\item \texttt{malloc}常维护空闲块信息，寻找合适空间，必要时拆分空闲块。
\item 没有合适空间时，可能向操作系统申请更多内存。
\item \texttt{free}可能把块放回分配器、合并相邻空闲块，供后续申请复用。
\item \texttt{realloc}可能在原处调整，也可能申请新块、复制数据再释放旧块。
\end{itemize}
\medskip\alert{这是帮助理解的实现思路，不是C标准规定的算法。}\par\medskip 频繁申请与释放可能带来fragmentation（碎片）；\texttt{free}也不保证进程占用立刻下降。
\note{补充页，按课堂时间决定是否展开，不实现malloc/calloc，不讲体系结构细节。}
\end{frame}

\section{Memory区域与对象的存储期}

\begin{frame}{先猜输出：每次进入函数，变量都会重新开始吗？}
\begin{Code}*{c}\lstinputlisting{code/lifetime.c}\end{Code}
\note{先留一分钟，让同学写下两行输出；不要先把static等同于全局可见。}
\end{frame}

\begin{frame}{答案：作用域与生命周期是两件事}
\begin{itemize}
\item 第一次：\texttt{1 1 11}；第二次：\texttt{1 2 12}。
\item \texttt{local}：每次进入函数体，开始新的automatic对象生命周期。
\item \texttt{count}：具有static storage duration，整个程序执行期间都存在。
\item \texttt{global}：同样具有static storage duration，值不会在函数返回时重置。
\end{itemize}
\medskip\alert{count只能在函数体内用名字访问，但它不会随着函数返回而消失。}
\end{frame}

\begin{frame}{Scope、storage duration、lifetime}
\begin{itemize}
\item \textbf{Scope（作用域）}：源码中，哪些位置能用这个名字？
\item \textbf{Storage duration（存储期）}：对象的存储存在多长时间？
\item \textbf{Lifetime（生命周期）}：对象在执行中何时存在、何时可以访问？
\end{itemize}
\medskip
\begin{tabular}{lll}
\textbf{对象} & \textbf{典型存储期} & \textbf{何时结束}\\[4pt]
普通块内局部变量 & automatic & 离开对应的块\\
全局变量／\texttt{static}局部变量 & static & 程序结束\\
\texttt{malloc}分配的对象 & allocated & 被释放或被成功重新分配
\end{tabular}

\note{本节限定普通非VLA示例，不引入thread storage duration。指针也是对象，拥有自己的生命周期。}
\end{frame}

\begin{frame}{Memory layout：常见Linux进程的示意图}
\begin{columns}[c]
\begin{column}{0.59\textwidth}
\centering
\begin{tikzpicture}[x=1cm,y=1cm,>=Stealth,
    band/.style={draw=VSPPrimary!65,line width=0.55pt},
    name/.style={anchor=west,font=\small\ttfamily\bfseries,text=VSPPrimary},
    desc/.style={anchor=west,font=\small,text=black!80}]
    % One address space, with a fixed width for every region.
    \fill[VSPPrimary!7] (0,0) rectangle (6,-0.72);
    \fill[black!2] (0,-0.72) rectangle (6,-2.22);
    \fill[VSPPrimary!7] (0,-2.22) rectangle (6,-2.94);
    \fill[VSPPrimary!9] (0,-2.94) rectangle (6,-4.22);
    \fill[white] (0,-4.22) rectangle (6,-4.86);
    \draw[band] (0,0) rectangle (6,-4.86);
    \foreach \y in {-0.72,-2.22,-2.94,-3.58,-4.22}
        \draw[band] (0,\y) -- (6,\y);
    % Read-only objects and instructions have a neutral background.
    \fill[white] (0.01,-4.22) rectangle (5.99,-4.85);
    \draw[band] (0,-4.86) -- (6,-4.86);
    \node[name] at (0.24,-0.36) {Stack};
    \node[desc] at (1.85,-0.36) {调用栈、部分局部对象};
    \node[name] at (0.24,-2.58) {Heap};
    \node[desc] at (1.85,-2.58) {动态分配区域};
    \node[name] at (0.24,-3.26) {.bss};
    \node[desc] at (1.85,-3.26) {零初始化静态对象};
    \node[name] at (0.24,-3.90) {.data};
    \node[desc] at (1.85,-3.90) {非零初始化静态对象};
    \node[name] at (0.24,-4.54) {.rodata};
    \node[desc] at (1.85,-4.54) {字符串字面量等只读数据};
    % Group the usual writable static-storage sections without adding a section.
    \draw[VSPPrimary,line width=1pt] (0,-2.94) rectangle (6,-4.22);
    \draw[VSPPrimary,line width=0.9pt] (6.08,-2.94) -- (6.24,-2.94)
        -- (6.24,-4.22) -- (6.08,-4.22);
    \node[rotate=90,font=\scriptsize\bfseries,text=VSPPrimary] at (6.48,-3.58) {Static};
    % Code sits in its own equal-width region below the data regions.
    \fill[white] (0,-4.86) rectangle (6,-5.50);
    \draw[band] (0,-4.86) rectangle (6,-5.50);
    \node[name] at (0.24,-5.18) {.text};
    \node[desc] at (1.85,-5.18) {机器指令};
    % RAM/ROM are teaching analogies for writable objects vs read-only content.
    \draw[VSPPrimary,line width=1pt] (6.68,0) -- (6.84,0)
        -- (6.84,-4.22) -- (6.68,-4.22);
    \node[rotate=90,font=\scriptsize\bfseries,text=VSPPrimary] at (7.08,-2.11) {RAM（类比）};
    \draw[black!65,line width=1pt] (6.68,-4.22) -- (6.84,-4.22)
        -- (6.84,-5.50) -- (6.68,-5.50);
    \node[rotate=90,font=\scriptsize\bfseries,text=black!65] at (7.08,-4.86) {ROM（类比）};
    % The gap is not a segment; arrows connect growth to the right region.
    \draw[->,VSPPrimary,very thick] (0.55,-0.87) -- (0.55,-1.55);
    \node[anchor=west,font=\scriptsize,text=VSPPrimary] at (0.77,-1.12) {Stack向低地址增长};
    \draw[->,VSPPrimary,very thick] (5.45,-2.07) -- (5.45,-1.39);
    \node[anchor=east,font=\scriptsize,text=VSPPrimary] at (5.23,-1.83) {Heap向高地址增长};
    \node[font=\small,text=black!55,fill=black!2,inner sep=2pt] at (3,-1.47) {空闲空间／其他映射};
    % Address direction is independent from allocation growth.
    \draw[->,black!45,line width=0.7pt] (-0.42,-5.35) -- (-0.42,-0.15);
    \node[anchor=east,font=\scriptsize,text=black!65,align=center] at (-0.52,-0.30) {高\\地址};
    \node[anchor=east,font=\scriptsize,text=black!65,align=center] at (-0.52,-5.15) {低\\地址};
\end{tikzpicture}
\end{column}
\begin{column}{0.37\textwidth}
\textbf{再按读写用途分组}
\medskip
\begin{itemize}
    \item \textbf{RAM类比}：Stack、Heap及Static区中的可写对象。
    \item \textbf{ROM类比}：.rodata的只读数据与.text的机器指令。
    \item \textbf{Static区}：将.bss与.data归组，对象存储期贯穿程序。
\end{itemize}
\medskip
\small\alert{Linux中.text／.rodata通常也驻留在物理RAM中；图中ROM仅类比只读属性。}
\par\medskip
\small 地址由下向上增大；箭头表示常见增长方向。C标准不保证这些排列与方向。
\end{column}
\end{columns}
\note{RAM与ROM在本图中仅作按访问属性的教学类比，并非物理内存或硬件地址区间。实际Linux的.text通常可读可执行、.rodata通常只读，二者均可驻留物理RAM，权限由映射控制。空闲空间与其他映射不据此统一规定读写权限。Static是这里为教学使用的分组标签，不是另一个ELF section；.bss与.data分别保留。具有static存储期的只读对象也可能放在.rodata，故不把此分组说成所有static对象唯一所在的区域。}
\end{frame}

\begin{frame}{这些对象都在哪？没写初值时会怎样？}
\begin{columns}[T]
\begin{column}{0.45\textwidth}
\begin{Code}*{c}\lstinputlisting{code/locations.inc}\end{Code}
\end{column}
\begin{column}{0.51\textwidth}
\scriptsize
\setlength{\tabcolsep}{4pt}
\renewcommand{\arraystretch}{1.3}
\begin{tabular}{lll}
\textbf{对象／数据} & \textbf{初始值} & \textbf{通常所在区域}\\
\texttt{global\_uninit} & 0 & .bss（Static）\\
\texttt{global\_init} & 7 & .data（Static）\\
\texttt{static\_global} & 8 & .data（Static）\\
\texttt{static\_local} & 9 & .data（Static）\\
\texttt{static\_zero} & 0 & .bss（Static）\\
\texttt{local} & 不确定值 & Stack／寄存器\\
\texttt{p} & 字面量的地址 & Stack／寄存器\\
\texttt{"hello"} & 字符及结尾零 & .rodata
\end{tabular}
\end{column}
\end{columns}
\medskip
\begin{itemize}
    \item \textbf{默认初始化行为，看存储期：}没写初值的全局／static对象会\textbf{零初始化}；普通局部\texttt{int local;}不会自动置零，\alert{先赋值，再读取}。
    \item 非零初值通常进入.data；显式\texttt{= 0}也通常进入.bss。static局部对象的生命周期贯穿程序。
    \item \alert{p保存地址，"hello"是所指数据；字符串字面量不能修改。}
\end{itemize}
\note{这里的“默认初始化行为”是在C17语境中解释未显式初始化时会发生什么，不套用C++的default-initialization分类。普通automatic局部对象未初始化时具有不确定值，不以“随机数”或“垃圾值”描述，更不能读取后猜测结果。位置是常见实现，对象可能被优化；字符串字面量具有static存储期。}
\end{frame}

\begin{frame}{.bss与.data：为什么两个程序大小差这么多？}
\begin{columns}[T]\begin{column}{0.53\textwidth}\begin{Code}*{c}\lstinputlisting{code/bss.c}\end{Code}\begin{Code}*{c}\lstinputlisting{code/data.c}\end{Code}\end{column}\begin{column}{0.43\textwidth}\begin{itemize}
\item 两组数组在运行时都需要存储。
\item 第一组默认全为0；第二组前三项非零，其余仍为0。
\item .bss通常只记录需要多少零初始化空间。
\item .data需要在文件中保存初始数据。
\end{itemize}
\end{column}\end{columns}\medskip\alert{可执行文件小，不代表程序运行时不占内存。}
\end{frame}

\begin{frame}{现场看一眼：size与objdump}
在\texttt{r5/}目录执行：\begin{Code}*{bash}\lstinputlisting{code/layout.sh}\end{Code}\begin{itemize}
\item \texttt{size}：比较text、data、bss所占空间。
\item \texttt{objdump -t}：观察符号\texttt{array}属于哪个section。
\item \texttt{int a;}与\texttt{int b = 0;}通常都进入.bss。
\end{itemize}

\note{工具演示约3分钟，只观察section归属。零初始化不是未初始化；具体文件大小取决于平台与工具链。}
\end{frame}

\begin{frame}{Stack frame：函数返回后发生了什么？}

\begin{columns}[c]
\begin{column}{0.43\textwidth}
\begin{tikzpicture}[node distance=0.12cm,every node/.style={draw=VSPPrimary,minimum width=4.6cm,minimum height=0.85cm,align=center}]
\node (c) {\texttt{function\_c}的frame};
\node[below=of c] (b) {\texttt{function\_b}的frame};
\node[below=of b] {\texttt{main}的frame};
\end{tikzpicture}
\end{column}
\begin{column}{0.52\textwidth}
\begin{itemize}
\item 函数调用通常建立stack frame（栈帧）。
\item 栈帧可能保存局部变量、返回地址与部分参数。
\item 参数也可能通过寄存器传递。
\item 函数返回后，普通局部对象生命周期结束。
\end{itemize}
\end{column}
\end{columns}
\medskip\alert{内存里的旧字节暂时没被覆盖，也不代表原来的对象还活着。}

\end{frame}

\section{谁负责free：指针、别名与所有权}

\begin{frame}{free：释放的是分配，不是指针变量}
\begin{columns}[c]
\begin{column}{0.46\textwidth}
\begin{Code}*{c}\lstinputlisting{code/free.inc}\end{Code}
\end{column}
\begin{column}{0.50\textwidth}
\centering
\begin{tikzpicture}[>=Stealth,font=\scriptsize]
    \node[font=\bfseries,text=black!70] at (0,1.15) {1. \texttt{free(p);} 之后};
    \node[draw=VSPPrimary,fill=VSPPrimary!6,minimum width=1.5cm,minimum height=0.65cm] (p1) at (-1.3,0.45) {\texttt{p}（旧地址）};
    \node[draw=black!45,dashed,fill=black!5,minimum width=2.2cm,minimum height=0.65cm,text=black!55] (b1) at (1.5,0.45) {已归还Heap（禁访）};
    \draw[->,dashed,VSPPrimary,thick] (p1.east) -- node[above,font=\tiny,text=VSPPrimary] {悬空!} (b1.west);
    \node[font=\bfseries,text=black!70] at (0,-0.35) {2. \texttt{p = NULL;} 之后};
    \node[draw=VSPPrimary,fill=VSPPrimary!6,minimum width=1.5cm,minimum height=0.65cm] (p2) at (-1.3,-1.05) {\texttt{p = NULL}};
    \node[draw=black!40,minimum width=1.2cm,minimum height=0.55cm,text=black!60] (n2) at (1.2,-1.05) {\texttt{NULL}};
    \draw[->,thick,black!55] (p2.east) -- (n2.west);
\end{tikzpicture}
\end{column}
\end{columns}
\medskip
\begin{itemize}
\item \texttt{free}结束该块的分配，但\textbf{不会自动把\texttt{p}改为\texttt{NULL}}。
\item \texttt{free(NULL)}安全无操作；主动置\texttt{p = NULL}可避免误用悬空指针。
\end{itemize}
\end{frame}

\begin{frame}{哪几种地址能传给free？}
\begin{columns}[c]
\begin{column}{0.48\textwidth}
\begin{Code}*{c}\lstinputlisting{code/invalid-free.inc}\end{Code}
\end{column}
\begin{column}{0.48\textwidth}
\centering
\begin{tikzpicture}[>=Stealth,font=\scriptsize]
    \node[font=\bfseries,text=black!75] at (1.65,1.25) {整租一套房，退房必须交整套房钥匙};
    \filldraw[fill=VSPPrimary!12,draw=VSPPrimary] (0,0) rectangle (1.1,0.65);
    \node at (0.55,0.325) {\texttt{p[0]}};
    \filldraw[fill=VSPPrimary!12,draw=VSPPrimary] (1.1,0) rectangle (2.2,0.65);
    \node at (1.65,0.325) {\texttt{p[1]}};
    \filldraw[fill=VSPPrimary!12,draw=VSPPrimary] (2.2,0) rectangle (3.3,0.65);
    \node at (2.75,0.325) {\texttt{p[2]}};
    \node[draw=VSPPrimary,fill=white,inner sep=2pt] (p0) at (0.55,-0.75) {\checkmark\ \texttt{free(p)}（块首地址）};
    \draw[->,thick,VSPPrimary] (p0.north) -- (0.55,0);
    \node[draw=black!50,dashed,fill=black!4,inner sep=2pt,text=VSPPrimary] (p1) at (2.35,-1.55) {$\times$\ \texttt{free(p + 1)}（内部地址）};
    \draw[->,thick,dashed,VSPPrimary] (p1.north) -- (1.65,0);
\end{tikzpicture}
\end{column}
\end{columns}
\medskip
\alert{只能传NULL，或与仍有效的动态分配返回值完全匹配的块首指针。}\par\smallskip\small 局部变量（Stack）、全局变量（Static）、块内部偏移地址（\texttt{p + 1}）和已释放的地址都不能传给\texttt{free}。
\end{frame}

\begin{frame}{两个指针指向同一块分配：p = NULL后，q安全吗？}
\begin{columns}[c]
\begin{column}{0.46\textwidth}
\begin{Code}*{c}\lstinputlisting{code/alias.inc}\end{Code}
\end{column}
\begin{column}{0.50\textwidth}
\centering
\begin{tikzpicture}[>=Stealth,font=\scriptsize]
    \node[draw=VSPPrimary,fill=VSPPrimary!6,minimum width=1.5cm,minimum height=0.65cm] (p) at (0,0.65) {\texttt{p = NULL}};
    \node[draw=VSPPrimary,fill=VSPPrimary!6,minimum width=1.5cm,minimum height=0.65cm] (q) at (0,-0.65) {\texttt{q}（旧地址）};
    \node[draw=black!45,dashed,fill=black!5,minimum width=2.4cm,minimum height=1.35cm,align=center,text=black!55] (x) at (3.6,0) {同一块动态分配\\（已被\texttt{free(p)}释放）};
    \draw[->,thick,dashed,VSPPrimary] (q.east) -- node[below=2pt,font=\tiny,text=VSPPrimary] {悬空!} (x.west);
\end{tikzpicture}
\medskip
\begin{itemize}\footnotesize
\item \texttt{int *q = p;}只复制地址，没有复制对象。
\item \texttt{p = NULL}不会自动更新别名\texttt{q}。
\item \alert{释放次数看分配次数，不看有几个指针。}
\end{itemize}
\end{column}
\end{columns}
\end{frame}

\begin{frame}{Memory leak：指针没了，堆上的分配还在}
\begin{columns}[c]
\begin{column}{0.48\textwidth}
\begin{Code}*{c}\lstinputlisting{code/leak.inc}\end{Code}
\end{column}
\begin{column}{0.48\textwidth}
\centering
\begin{tikzpicture}[>=Stealth,font=\scriptsize]
    \node[draw=black!45,dashed,fill=black!4,minimum width=2.1cm,minimum height=1.1cm,align=center,text=black!55] (stack) at (0,0) {\texttt{f}的栈帧已销毁\\局部指针\texttt{p}消失};
    \node[draw=VSPPrimary,fill=VSPPrimary!12,minimum width=2.3cm,minimum height=1.1cm,align=center] (heap) at (3.9,0) {Heap中的100个\texttt{int}\\仍然占着内存};
    \draw[->,dashed,black!35,thick] (stack.east) -- node[above,text=VSPPrimary,font=\bfseries] {$\times$ 断联} (heap.west);
\end{tikzpicture}
\end{column}
\end{columns}
\medskip
\begin{itemize}
\item \texttt{p}是automatic局部指针，离开函数\texttt{f}后自动销毁。
\item 但Heap里的对象不会自动释放；丢失唯一地址后，程序再也无法\texttt{free}它。
\end{itemize}
\alert{“指针生命周期结束”不等于“所指对象自动释放”。}
\end{frame}

\begin{frame}{回看create\_point：约定谁负责free}
\begin{Code}*{c}\lstinputlisting[firstline=12,lastline=22,firstnumber=12]{code/create-point.c}\end{Code}
\begin{itemize}
    \item \texttt{create\_point}把动态对象交给调用者（转移\textbf{ownership}）；返回\texttt{NULL}表示失败。
    \item \texttt{main}接管返回的指针，使用结束后负责\texttt{free(p)}。
    \item 局部结构体\texttt{new\_point}随作用域自动结束生命周期，不能\texttt{free(\&new\_point)}。
\end{itemize}
\end{frame}

\begin{frame}{链表：每个节点都是一次独立分配}
\begin{Code}*{c}\lstinputlisting{code/list-node.inc}\end{Code}
\begin{center}
\begin{tikzpicture}[>=Stealth,node distance=0.7cm,
every node/.style={draw=VSPPrimary,minimum height=0.8cm,inner sep=7pt}]
\node (head) {head};
\node[right=of head] (a) {20 | next};
\node[right=of a] (b) {10 | next};
\node[right=of b] (nil) {NULL};
\draw[->,thick] (head)--(a);
\draw[->,thick] (a)--(b);
\draw[->,thick] (b)--(nil);
\end{tikzpicture}
\end{center}
\alert{释放head指向的节点，不会自动释放next指向的节点。}

\end{frame}

\begin{frame}{释放链表：为什么必须先保存next，再free？}
\begin{columns}[c]
\begin{column}{0.47\textwidth}
\begin{Code}*{c}\lstinputlisting{code/list-free.inc}\end{Code}
\end{column}
\begin{column}{0.49\textwidth}
\centering
\begin{tikzpicture}[>=Stealth,font=\scriptsize]
    \node[font=\bfseries,text=VSPPrimary] at (1.8,1.35) {$\times$ 错误：先\texttt{free(head)}，桥已经炸断};
    \node[draw=black!45,dashed,fill=black!6,minimum width=1.6cm,minimum height=0.6cm,text=black!50] (wa) at (0.1,0.6) {已\texttt{free}节点};
    \node[draw=VSPPrimary,fill=VSPPrimary!8,minimum width=1.5cm,minimum height=0.6cm] (wb) at (3.5,0.6) {下一节点};
    \draw[->,dashed,VSPPrimary,thick] (wa.east) -- node[above,font=\tiny,text=VSPPrimary] {读\texttt{next} = UAF!} (wb.west);

    \node[font=\bfseries,text=black!75] at (1.8,-0.15) {$\checkmark$ 正确：先用\texttt{next}探路，再拆当前桥};
    \node[draw=VSPPrimary,fill=VSPPrimary!8,minimum width=1.6cm,minimum height=0.6cm] (ca) at (0.1,-0.85) {\texttt{head}当前节点};
    \node[draw=VSPPrimary,fill=VSPPrimary!8,minimum width=1.5cm,minimum height=0.6cm] (cb) at (3.5,-0.85) {下一节点};
    \draw[->,thick,VSPPrimary] (ca.east) -- (cb.west);
    \node[draw=VSPPrimary,fill=white,inner sep=2pt] (nx) at (3.5,-1.7) {1. \texttt{next}先指向它};
    \draw[->,thick,VSPPrimary] (nx.north) -- (cb.south);
    \node[font=\tiny,text=VSPPrimary] at (0.1,-1.5) {2. 安全\texttt{free(head)}};
\end{tikzpicture}
\end{column}
\end{columns}
\medskip
\alert{先保存next指针，再free当前节点；否则读取已释放节点的next就是use-after-free。}
\end{frame}

\begin{frame}{内存错误小结：先问哪条规则被破坏了}
\small
\setlength{\tabcolsep}{5pt}
\renewcommand{\arraystretch}{1.28}
\begin{tabular}{lll}
\textbf{错误类型} & \textbf{破坏了哪条规则} & \textbf{本节典型现场}\\[3pt]
Out-of-bounds（越界） & 下标超出有效分配范围 & 申请\texttt{n}个元素却访问\texttt{a[n]}\\
Uninitialized read & 读取前未写入确定值 & 局部\texttt{int local;}或\texttt{malloc}后直接读\\
Dangling / UAF & 所指对象生命周期已结束 & 返回\texttt{\&point}、先\texttt{free(head)}再读\texttt{next}\\
Double free（重复释放） & 同一块分配被\texttt{free}两次 & 别名\texttt{free(p)}后又执行\texttt{free(q)}\\
Memory leak（内存泄漏） & 丢失地址，遗漏释放路径 & \texttt{p = realloc(p, ...)}失败覆盖原指针
\end{tabular}
\medskip\par\alert{这些问题没有按stack、heap、global划分的固定严重程度。}
\par\smallskip 越界、use-after-free、double free等问题，可结合\textbf{上节课的AddressSanitizer}检查。
\end{frame}

\begin{frame}{带走这四个问题}
\begin{itemize}
\item \textbf{在哪里？}区分指针变量与所指对象；布局是实现模型。
\item \textbf{活多久？}区分scope、automatic/static/allocated存储期。
\item \textbf{能访问吗？}检查初始化、边界、NULL与生命周期。
\item \textbf{谁释放？}每次分配都有明确责任；失败路径也要清理。
\end{itemize}
\medskip\alert{地址还在，不代表对象还在；函数返回，也不代表动态对象被释放。}
\end{frame}
\VSPendframe{}
\appendix
\section{课后练习：指针与宏}

\begin{frame}{课后练习一：判断对错，并说明理由}
\begin{itemize}
\item A. 普通局部变量一定活到整个程序结束。
\item B. 动态分配的对象可以跨越一次函数调用继续存在。
\item C. 给指针赋值\texttt{NULL}会自动释放它原先指向的块。
\item D. 全局变量可以通过\texttt{free(\&global)}释放。
\item E. 对一个仍有效的分配，多个指针可以同时指向它。
\end{itemize}
\medskip\alert{先独立判断，再用“对象的生命周期”解释。}
\end{frame}

\begin{frame}{课后练习一答案：B、E正确}
\begin{itemize}
\item A错：普通automatic局部对象在离开对应块时结束生命周期。
\item B对：动态分配对象的存在不依赖创建它的函数是否已经返回。
\item C错：赋值改变的是指针变量；原分配可能因此泄漏。
\item D错：全局对象不来自动态分配，不能这样释放。
\item E对：多个别名共享同一个对象，但必须协调访问与释放责任。
\end{itemize}

\end{frame}

\begin{frame}{课后练习二：移动形参，会改变调用者吗？}
\begin{Code}*{c}\lstinputlisting{code/pointer.c}\end{Code}\alert{输出是4 4，还是4 6？画出调用前后的两个p。}
\end{frame}

\begin{frame}{课后练习二答案：输出4 4}
\begin{itemize}
\item C函数参数按值传递，\texttt{move}得到地址值的副本。
\item 函数里的\texttt{p = p + 1}只改变形参，不改变调用者的\texttt{p}。
\item \texttt{*a}仍然是\texttt{a[0]}；调用者的\texttt{p}也仍指向\texttt{a[0]}。
\end{itemize}
\begin{Code}*{c}\lstinputlisting{code/pointer-change.inc}\end{Code}\alert{改对象内容传对象地址；改调用者的指针，就传那个指针的地址。}
\end{frame}

\begin{frame}{课后练习三：加2，同时返回原值}
补全函数：调用后\texttt{x}加2，返回修改前的值。\begin{Code}*{c}\lstinputlisting{code/add-two-question.inc}\end{Code}\begin{itemize}
\item 调用者希望自己的\texttt{x}被修改，需要传什么？
\item 如何保存修改前的值？
\item 前提：\texttt{x}指向有效\texttt{int}对象，原值加2不发生有符号溢出。
\end{itemize}

\end{frame}

\begin{frame}{课后练习三答案：先保存，再修改}
\begin{Code}*{c}\lstinputlisting{code/add-two.c}\end{Code}输出\texttt{0 2}。\alert{指针让函数找到调用者的对象；它本身仍然是按值传递。}
\end{frame}

\begin{frame}{课后练习四：宏不会替你补括号}
\begin{Code}*{c}\lstinputlisting{code/macro-mul-question.inc}\end{Code}\begin{itemize}
\item 先做macro expansion（宏展开），再按运算符优先级计算。
\item 展开后是\texttt{16 / 1 + 1 * 3 - 1}。
\item 它和\texttt{16 / ((1 + 1) * (3 - 1))}一样吗？
\item 怎么修改\texttt{MUL}，让这一次计算得到期望结果？
\end{itemize}

\end{frame}

\begin{frame}{课后练习四答案：原式为18，修复后为4}
\begin{Code}*{c}\lstinputlisting{code/macro-mul.c}\end{Code}\begin{itemize}
\item 参数各自加括号，保证每个实参被当作完整表达式。
\item 整体再加括号，保证整个展开结果参与外围运算。
\end{itemize}

\end{frame}

\begin{frame}{课后练习五：用条件运算符写MIN}
\begin{Code}*{c}\lstinputlisting{code/macro-min-question.inc}\end{Code}\begin{itemize}
\item 比较\texttt{a}与\texttt{b}，较小时选择哪一个？
\item 每个参数和整个表达式都要考虑括号。
\item 追问：\texttt{MIN(i++, j)}会让\texttt{i}增加几次？
\end{itemize}

\end{frame}

\begin{frame}{课后练习五答案：括号解决不了重复求值}
\begin{Code}*{c}\lstinputlisting{code/macro-min.c}\end{Code}这里先比较\texttt{1 < 5}，再求真分支\texttt{i++}，输出\texttt{2 3}。\par\alert{避免把带副作用的表达式传给这类宏；普通整数比较可用函数。}
\note{这里条件运算符的条件与选中分支之间有顺序保证，该例不是未定义行为。不要用MIN(i++, i++)作为对比，它会引入未定序修改问题。}
\end{frame}

\begin{frame}{课后练习六：打印类型名和sizeof结果}
写一个宏\texttt{PS(type)}，例如\texttt{PS(long long)}。\begin{Code}*{c}\lstinputlisting{code/macro-size-question.inc}\end{Code}\begin{itemize}
\item \texttt{sizeof}的结果类型是\texttt{size\_t}，使用\texttt{\%zu}。
\item 如何把宏参数\texttt{long long}变成字符串？提示：\texttt{\#type}。
\item \alert{示例里的8只是常见结果，不是所有平台的固定答案。}
\end{itemize}

\end{frame}

\begin{frame}{课后练习六答案：\#把参数字符串化}
\begin{Code}*{c}\lstinputlisting{code/macro-size.c}\end{Code}\begin{itemize}
\item \texttt{\#type}是stringification（字符串化）。
\item 相邻字符串字面量在翻译时会拼接。
\item \texttt{PS(int);}相当于打印\texttt{"Size of " "int" ": \%zu\\n"}并传入\texttt{sizeof(int)}。
\end{itemize}

\end{frame}

\begin{frame}{课后综合题：这段代码有哪些问题？}
\begin{Code}*{c}\lstinputlisting{code/review-question.inc}\end{Code}假定已包含头文件，代码位于\texttt{main}中。\par\medskip\alert{逐行分析；不要尝试预测未定义行为后的输出。}
\end{frame}

\begin{frame}{课后综合题答案：沿着分配与生命周期找问题}
\begin{itemize}
\item 没有检查\texttt{malloc}是否返回\texttt{NULL}。
\item 只申请两个\texttt{int}，合法下标是0和1；\texttt{p[2]}越界。
\item \texttt{free(p)}之后\texttt{q}也失去有效所指对象；读取\texttt{q[0]}是use-after-free。
\item 最后\texttt{free(q)}重复释放同一分配。
\item \texttt{p = NULL}只更新\texttt{p}，无法修复\texttt{q}。
\end{itemize}
\medskip\alert{正确顺序：检查分配 → 初始化 → 范围内访问 → 统一释放 → 停止使用旧别名。}
\end{frame}

\end{document}
